\paragraph{Thí nghiệm và kết quả}
\subparagraph{Fisher ratio}

Để đánh giá mức độ phân biệt của từng đặc trưng, ta sử dụng tiêu chí Fisher
ratio. Giá trị Fisher ratio càng lớn cho thấy đặc trưng đó có khả năng phân
tách các lớp càng tốt.

\tablename~\ref{tab:fisher-top-features} trình bày 5 đặc trưng có giá trị
Fisher ratio cao nhất.
\begin{table}[htbp]
  \centering
  \caption{Top 5 đặc trưng theo Fisher ratio}
  \label{tab:fisher-top-features}
  \import{\tablespath}{fisher-top-features}
\end{table}

Từ kết quả trên, có thể thấy rằng đặc trưng \textit{Elevation} có giá trị
Fisher ratio cao nhất, cho thấy đây là biến quan trọng nhất trong việc phân
biệt các lớp. Điều này gợi ý rằng độ cao có ảnh hưởng đáng kể đến sự khác biệt
giữa các nhóm dữ liệu.

Ngoài ra, các đặc trưng như \textit{Wilderness\_Area\_3} cũng có giá trị tương
đối cao, cho thấy vai trò đáng kể trong việc phân tách lớp. Ngược lại, các đặc
trưng thuộc nhóm \textit{Soil Type} có giá trị Fisher ratio thấp hơn, cho thấy
khả năng phân biệt yếu hơn khi xét riêng lẻ.

\subparagraph{Nhận xét}
Kết quả này cho thấy rằng chỉ một số ít đặc trưng đóng vai trò quan trọng
trong việc phân loại, trong khi nhiều đặc trưng khác có đóng góp hạn chế. Điều
này gợi ý rằng có thể áp dụng các phương pháp lựa chọn đặc trưng (feature
selection) để giảm số chiều của dữ liệu mà vẫn giữ được hiệu năng của mô hình.

Tuy nhiên, cần lưu ý rằng Fisher ratio đánh giá từng đặc trưng một cách độc
lập, do đó không phản ánh được mối quan hệ hoặc tương tác giữa các đặc trưng.
Vì vậy, việc kết hợp với các phương pháp khác (như Lasso hoặc các mô hình học
máy) có thể mang lại cái nhìn toàn diện hơn.

\subparagraph{Đường biên quyết định của LDA}

\figurename~\ref{fig:lda-decision-boundary} minh họa đường biên quyết định của
mô hình LDA sau khi chiếu dữ liệu xuống không gian hai chiều. Có thể thấy rằng
các biên quyết định của LDA là các đường thẳng, phù hợp với giả định rằng các
lớp chia sẻ cùng ma trận hiệp phương sai. Điều này dẫn đến việc không gian đặc
trưng được chia thành các vùng tuyến tính, mỗi vùng tương ứng với một lớp.
Trong hình, các lớp dữ liệu có xu hướng phân bố dọc theo một cấu trúc gần
tuyến tính, do đó LDA vẫn có khả năng phân tách tương đối tốt ở một số vùng.
Tuy nhiên, ở các khu vực mà các lớp chồng lấn mạnh hoặc có cấu trúc phi tuyến,
các biên tuyến tính của LDA không đủ linh hoạt để phân tách hoàn toàn, dẫn đến
sự xen lẫn giữa các lớp. Điều này phản ánh hạn chế của LDA là mô hình có độ
linh hoạt thấp (high bias), nên mặc dù ổn định và tổng quát hóa tốt, nhưng có
thể không đạt hiệu năng tối ưu khi dữ liệu có ranh giới phân lớp phức tạp.

\begin{figure}[htbp]
  \centering
  \includegraphics[width=0.9\textwidth]{lda-2d-decision-boundary}
  \caption{Đường biên quyết định của LDA trong không gian hai chiều}
  \label{fig:lda-decision-boundary}
\end{figure}
